% Target venue: Clinical Psychology Review (Elsevier), elsarticle class.
% Fallbacks: Perspectives on Psychological Science (SAGE/APA 7), Psychological Medicine (CUP, 4,500 words).
% Build: latexmk -pdf main.tex   (or: tectonic main.tex)
\documentclass[preprint,12pt,authoryear]{elsarticle}
\usepackage[utf8]{inputenc}
\usepackage{amsmath,booktabs,array}
\usepackage[hidelinks]{hyperref}
\biboptions{authoryear,round}
\journal{Clinical Psychology Review}

\usepackage{xcolor}
\newcommand{\citegap}[1]{\textcolor{red}{[CITE: #1]}}

\begin{document}
\begin{frontmatter}
\title{Why Does CBT Fail? A Critical Meta-Review of Non-Response, Weak Mechanism Evidence and Delivery Limits, with Executive Function and Adult ADHD as a Test Case}

\author{Brando Miranda\corref{cor1}}
\ead{brando9@stanford.edu}
\cortext[cor1]{Corresponding author}
\address{Stanford University, Stanford, CA, USA}

\begin{abstract}
Cognitive behavioral therapy (CBT) is the best-studied psychotherapy for anxiety and depression, yet only about half of treated patients respond, and ``CBT does not work'' is used to mean very different things. We map ten reasons given for why CBT fails and grade each against meta-analytic evidence: (R1) the cognitive skill demands executive function that anxiety and attention-deficit/hyperactivity disorder (ADHD) degrade; (R2) cognitive restructuring is not the active ingredient; (R3) headline effects are inflated by waitlists, bias and measurement; (R4--R6) non-response, dropout, relapse and under-recorded harm; (R7--R8) patient and context factors, including neurodivergence, comorbidity, culture and access; (R9) therapist and delivery factors; and (R10) an under-tested theory of change. About 440 sources were DOI-verified. Dropout (about 16\% before and 26\% during CBT), non-response (fewer than half meeting response or remission criteria in 92 psychotherapy studies), relapse (29\% within one year, 54\% within two after cognitive therapy ends), biased and heterogeneously measured trials and under-reported harms are supported, but almost all are shared with other treatments. The cognitive mechanism is under-tested, and component studies repeatedly fail to find added benefit from restructuring, though they are underpowered. The executive-function argument, our deepest case study, survives only in a narrow form: reappraisal draws on working memory in the laboratory, evidence that stress abolishes it is contested, and we found no study of ADHD as a moderator of CBT for anxiety in adults. CBT keeps a moderate average effect under almost every sensitivity analysis; the review locates where and why it falls short and closes with preregistrable predictions. \textbf{Draft 3; for the author's review.}
\end{abstract}

\begin{keyword}
cognitive behavioral therapy \sep cognitive restructuring \sep executive function \sep anxiety \sep ADHD \sep non-response \sep dropout \sep relapse \sep therapist effects \sep critical review
\end{keyword}
\end{frontmatter}


\section{Introduction}
Cognitive behavioral therapy (CBT) is among the most heavily tested psychotherapies. A meta-analysis of 41 placebo-controlled trials in anxiety and related disorders (total $N=2{,}843$) found that CBT outperformed placebo on target symptoms with a moderate effect (Hedges' $g=0.56$) \citep{carpenter2018cbt}. The same literature also shows a ceiling: across 87 studies only about half of patients respond to CBT for anxiety disorders, and the figure depends on how response is defined \citep{loerinc2015response}; a larger meta-analysis of 100 articles reports a mean remission rate of 51.0\% \citep{springer2018remission}. A therapy can work on average and still leave open \emph{why} it works and \emph{for whom}.

The textbook account of CBT's cognitive component is a top-down loop. A situation triggers an automatic thought (``I am going to fail this talk''); the patient records it, names the distortion, weighs evidence for and against, and generates a more balanced thought; the feeling is expected to follow. This review first maps ten reasons given for why CBT fails (Section~\ref{sec:taxonomy}) and then grades the evidence for each. Its deepest case study is a narrow question about that loop: \emph{does the mechanism ask the patient to use a capacity that anxiety, and in some patients attention-deficit/hyperactivity disorder (ADHD), has already degraded?} If so, the cognitive component should depend on baseline executive capacity, should be harder to use when stressed than when calm, and should contribute less to outcome than behavioural elements that do not demand effortful reasoning.

The argument has a popular version. Judson Brewer, a clinician-researcher, argues that anxiety takes the prefrontal cortex ``off-line,'' so that willpower and reasoning are unavailable at the moments when they are needed \citep{brewer2016habit,brewer2021unwinding}, and he has applied this directly to CBT \citep{brewer2026cbt}. Our aim is to find out how much of this survives contact with the primary literature. Our answer, in brief, is that a narrow version survives and the strong version does not. The same standard is applied to the other reasons: dropout and non-response (Section~\ref{sec:nonresponse}), patient and context factors (Section~\ref{sec:patients}), therapist and delivery factors (Section~\ref{sec:delivery}), and the theory, measurement and evidence-quality critiques (Section~\ref{sec:mechanism}).

\paragraph{What we find} (1) Reappraisal does draw on working memory in laboratory studies, but evidence that acute stress abolishes it is thin and contested. (2) The direct evidence that executive function moderates CBT outcome is sparse and mostly null for overall outcome, with the clearest signal at the level of the cognitive skill itself. (3) Component studies repeatedly find that adding cognitive restructuring to exposure or behavioural activation does not add measurable benefit, but these studies are underpowered and a minority find a benefit. (4) CBT's headline effects shrink with active controls and bias adjustment; whether they shrink over time is disputed. (5) In adult ADHD, CBT is modestly effective, mostly as an adjunct to medication, and the plausible active ingredient is behavioural compensation; we found no trial testing ADHD as a moderator of CBT for anxiety in adults. (6) Beyond that argument, dropout (about 16\% before and 26\% during CBT), non-response, relapse and under-recorded harm are real and quantified, but mostly shared with other psychological treatments; patient, context and delivery factors shift outcomes without isolating a CBT-specific failure; and the theory of change, the outcome measures and the trial base are the weakest links.

\paragraph{Contributions} (i) An explicit decomposition of cognitive restructuring into executive demands (Section~\ref{sec:demands}); (ii) an evidence-graded review of the stress-and-prefrontal premise, including a check of the ``off-line'' claim against its sources (Section~\ref{sec:stress}); (iii) a review of efficacy, bias and component evidence (Sections~\ref{sec:evidence}--\ref{sec:ingredients}); (iv) a review of alternatives and of CBT in adult ADHD (Sections~\ref{sec:alternatives}--\ref{sec:adhdcbt}); (v) the direct moderator evidence (Section~\ref{sec:moderators}); (vi) a graded map of ten reasons CBT fails (Section~\ref{sec:taxonomy}) and reviews of non-response, patient, delivery and mechanism factors (Sections~\ref{sec:nonresponse}--\ref{sec:mechanism}); (vii) preregistrable predictions (Section~\ref{sec:agenda}).

\section{Method}\label{sec:method}
This is a critical (theoretical) review, not a PRISMA-complete systematic review. It selects the literatures needed to evaluate one mechanistic argument and states what each can and cannot show. Searches ran in October 2026 across PubMed, Europe PMC, Crossref, publisher pages and the reference lists of key meta-analyses, in eight slices, the first four built for the executive-function argument: (a) efficacy, bias and head-to-head comparisons; (b) mechanisms and component (dismantling) studies, including exposure and third-wave alternatives; (c) executive function, stress and emotion regulation, including executive function as a moderator of treatment; and (d) ADHD (CBT, medication, mindfulness, executive function, anxiety comorbidity); and, added in Draft 3 to cover all the reasons CBT is said to fail, (e) non-response, dropout, relapse and harm; (f) patient and context factors (comorbidity, neurodivergence, culture, access); (g) therapist, delivery and real-world effectiveness; and (h) theory, mechanism and measurement critiques. Priority was given to meta-analyses, umbrella reviews, component trials and moderator studies, and to the strongest critiques and defences of CBT.

\paragraph{Verification} Every reference was resolved through Crossref (title, authors and year compared with the intended citation). Every quoted number was read from the abstract or full text, not recalled. Roughly 440 candidate sources were verified this way; several candidate references proposed from memory were dropped because the DOI did not resolve or resolved to a different paper. Sources for which only a partial abstract or a secondary summary was available are flagged in \texttt{deep-research/sources/}. Popular sources (books, talks, blogs) are cited only for the claims their authors make.

\paragraph{Limitations of the method} Single-reviewer screening; no pre-registered protocol; no formal risk-of-bias assessment of included reviews; evidence from depression is used where anxiety evidence on a question is thin, and flagged when so. A PRISMA-style umbrella search is the obvious next step (Section~\ref{sec:agenda}).

\section{A map of the reasons CBT fails}\label{sec:taxonomy}
``CBT does not work'' is not one claim. It bundles at least ten different ones, which differ in what they say fails (the average effect, the individual response, the mechanism, the evidence or the delivery), in how well they are supported, and in whether they are specific to CBT at all. Table~\ref{tab:reasons} lists them with the verdict this review reaches and the section that holds the evidence. We use four verdicts: \textbf{supported} (consistent evidence, quantified), \textbf{partly supported} (real but smaller or more conditional than the popular version), \textbf{contested} (credible evidence on both sides) and \textbf{open} (little or no direct evidence).

Two cautions apply throughout. First, many of the best-documented limits (dropout, relapse, comorbidity, harm under-reporting, therapist variability) are limits of psychological treatment in general, not of CBT in particular; we say so each time the source does. Second, a limit on certainty is not evidence of ineffectiveness: CBT keeps a moderate average effect under almost every sensitivity analysis we found \citep{plessen2023multiverse,carpenter2018cbt}. The map is therefore a list of \emph{where and why} CBT falls short of its reputation, not a verdict that it fails.

\begin{table}[ht]
\centering\small
\renewcommand{\arraystretch}{1.15}\begin{tabular}{@{}p{0.05\textwidth}>{\raggedright\arraybackslash}p{0.42\textwidth}>{\raggedright\arraybackslash}p{0.27\textwidth}>{\raggedright\arraybackslash}p{0.2\textwidth}@{}}
\toprule
 & Reason & Verdict & Where \\ \midrule
R1 & The cognitive skill demands executive function that anxiety (and ADHD) degrade & contested; direct moderator evidence open & Sections~\ref{sec:demands}--\ref{sec:adhd}, \ref{sec:moderators} \\
R2 & Cognitive restructuring is not the active ingredient & partly supported (underpowered trials) & Section~\ref{sec:ingredients} \\
R3 & Headline effects are inflated by waitlists, bias, low-quality trials and measurement choices & supported & Sections~\ref{sec:evidence}, \ref{sec:mechanism} \\
R4 & Many patients do not respond or drop out & supported (not CBT-specific) & Section~\ref{sec:nonresponse} \\
R5 & Gains erode after treatment ends & partly supported & Section~\ref{sec:nonresponse} \\
R6 & Harms and deterioration are under-recorded & supported for reporting; CBT lowers deterioration versus controls & Section~\ref{sec:nonresponse} \\
R7 & Patient factors (personality disorder, chronicity, adversity, medical comorbidity, neurodivergence) & partly supported; mostly not CBT-specific; autism evidence thin & Section~\ref{sec:patients} \\
R8 & Context: WEIRD trials, culture, waiting times, deprivation, access & supported & Section~\ref{sec:patients} \\
R9 & Delivery: therapist variability, drift, weak training, dose, format & partly supported; causal evidence missing & Section~\ref{sec:delivery} \\
R10 & The theory of change is under-tested; common factors may explain more & contested & Section~\ref{sec:mechanism} \\ \bottomrule
\end{tabular}
\vspace{4pt}\caption{Ten reasons given for why CBT fails, with the verdict of this review.}\label{tab:reasons}
\end{table}

The executive-function argument (R1) is the review's deepest case study, because it makes a testable, mechanism-level prediction and because it was the starting point of the project. It is one reason among ten, and Sections~\ref{sec:nonresponse}--\ref{sec:mechanism} treat the others at the same standard of evidence.

\section{What cognitive restructuring asks of the patient}\label{sec:demands}
Consider the standard thought-record exercise. The patient must (a) \emph{notice} an automatic thought while it is occurring, (b) hold the thought, the supporting evidence, the contradicting evidence and a candidate alternative in mind at once, (c) suppress the pull of the automatic thought while evaluating it, and (d) shift from being inside the feeling to observing it. Mapped onto the best-known model of executive function, these are three separable components: updating and holding information (working memory), inhibiting a prepotent response, and shifting between mental sets \citep{miyake2000unity}.

\paragraph{Laboratory evidence that reappraisal draws on executive resources} In event-related-potential studies, high working-memory load eliminated the effect of reappraisal on the late positive potential \citep{gan2017load,adamczyk2022load}. Working-memory capacity predicted reappraisal success \citep{schmeichel2008memory}, and reappraising late in an emotional response, once emotion is already intense, costs self-control resources \citep{sheppes2008costs}. People also choose differently as intensity rises: they prefer distraction to reappraisal in high-intensity negative situations and reappraisal in low-intensity ones \citep{sheppes2011choice}.

\paragraph{Evidence that it works, and can be trained, when load is low} Instructed reappraisal reduces negative experience in the laboratory \citep{hofmann2009anxiety} and can be trained \citep{denny2014training}. So the claim is not that reappraisal fails; it is that it \emph{costs} executive resources, which matters when those resources are scarce.

\paragraph{Two clarifications} First, this is a task analysis plus laboratory support, not a treatment finding. Clinicians scaffold the load (worksheets, a therapist who prompts for evidence, graded practice). Second, executive function is not a single resource; the argument concerns working memory, inhibition and shifting, not general intelligence. We return to both in Section~\ref{sec:objections}.

\section{Stress, anxiety and the prefrontal executive system}\label{sec:stress}
\subsection{Mechanistic neuroscience}
The most developed account is due to Arnsten. Her review of stress signalling describes the prefrontal cortex (PFC) as the most evolved brain region and also the one most sensitive to the detrimental effects of stress exposure \citep{arnsten2009stress}. The proposed mechanism is an inverted-U in catecholamine signalling: moderate norepinephrine and dopamine levels support prefrontal network firing and working memory, while the high levels of uncontrollable stress weaken them and strengthen amygdala and striatal habit circuitry. A later review summarises the result as stress rapidly flipping the brain from reflective to reflexive control of behaviour \citep{arnsten2015stress}. Acute uncontrollable stress impairs PFC-dependent working memory in monkeys and rodents \citep{arnsten1998noise}. In humans, a meta-analysis finds that acute stress impairs working memory and cognitive flexibility, with more nuanced effects on inhibition and many moderators \citep{shields2016stress}, and imaging suggests a reallocation of resources toward salience networks rather than a switch-off \citep{hermans2014networks}.

\paragraph{Where the ``off-line'' language comes from} The phrase traces to a hedged line in a 1998 monkey study (``Stress may take the PFC `off-line'''; \citealt{arnsten1998noise}), and its author still uses it in quotation marks for \emph{uncontrollable} stress \citep{arnsten2020guanfacine}. Brewer and Roy apply it to worry and cite Arnsten's 2009 and 2015 reviews, which are not studies of worry \citep{brewer2021habit}; their app-based trial ($n=65$, treatment-as-usual control) did not test a prefrontal mechanism \citep{roy2021app}. We therefore read ``offline'' as a popularisation of \emph{impaired, in a graded and dose-dependent way}, mainly under uncontrollable stress.

\subsection{Does stress abolish reappraisal in humans?}
Evidence for this specific link is contested. The founding study found that participants stressed before learning showed no fear reduction after regulation training \citep{raio2013stress}, and a review summarises that acute stress may impair the persistent inhibition of fear, presumably by altering prefrontal function \citep{raio2015influence}. Five later studies ($n=54$--$98$) found no group-level impairment \citep{shermohammed2017stress,sandner2021withstands,langer2021delayed,langer2023rapid,hamza2024acceptance}, though two reported sex- or cortisol-linked effects. Their laboratory stressors were moderate, so the null results may not transfer to clinical-level arousal; that gap is itself an open question.

\subsection{Cognitive psychology of anxiety}
Attentional Control Theory proposes that anxiety impairs the inhibition and shifting functions of executive control, so that anxious individuals are more captured by threat and worse at redirecting attention \citep{eysenck2007act}. The theory predicts that \emph{efficiency} (effort and time) suffers more than \emph{effectiveness} (accuracy), because anxious people compensate with extra effort. Meta-analytic support exists at group level \citep[$g=-0.58$;][]{shi2019attentional} and for working memory \citep[$g=-0.334$;][]{moran2016working}, but with qualifiers: a meta-analysis of 32 studies ($N=13{,}084$) found generalised anxiety disorder (GAD) associated with poorer cognitive flexibility and working memory but not inhibitory control, and with no association with any executive domain in older adults \citep{nguyen2025executive}. Working memory was unrelated to social anxiety in one study \citep{harrison2025executive}, and one study found \emph{enhanced} control over negative distractors in GAD \citep{hallion2019enhanced}.

\section{Executive function in ADHD}\label{sec:adhd}
ADHD is heterogeneous: executive difficulties are common but not universal. A meta-analysis of 83 studies found that groups with ADHD underperform controls on every executive task class examined, most strongly on response inhibition, vigilance, working memory and planning, with medium effect sizes ($0.46$--$0.69$) \citep{willcutt2005ef}. The same authors caution that executive weaknesses are neither necessary nor sufficient to cause all cases of ADHD, and fewer than half of the (mostly child) samples show impairment on any single executive test \citep{willcutt2005ef}. Measured more carefully, impairment is more prevalent: in a sample of 136 children, 89\% of those with ADHD were impaired on at least one executive function (62\% working memory, 27\% inhibitory control, 38\% set shifting) \citep{kofler2019heterogeneity}.

Anxiety is common in adult ADHD: a meta-analysis of 58 studies ($N=18{,}821$) found higher anxiety scores in adults with ADHD than without ($g=0.77$; $g=0.95$ against non-clinical and $g=0.44$ against clinical controls) \citep{hoskins2026anxiety}. Whether anxiety adds executive load in ADHD is less clear: in a sample of 339 children, trait anxiety did not uniquely predict working memory or inhibitory control, nor moderate the ADHD effect \citep{marsh2024anxiety}. The stacking assumption (low baseline executive capacity, lowered further by anxiety) is therefore plausible but, in children at least, not supported by the one careful test we found.

\section{What the efficacy evidence shows}\label{sec:evidence}
\subsection{The case for CBT}
Against placebo, CBT outperforms on anxiety target symptoms ($g=0.56$ across 41 trials; \citealt{carpenter2018cbt}). Against inactive controls, effects are larger: a meta-analysis of 144 trials found $g=0.75$ (depression) to $0.88$ (social anxiety disorder), with $0.80$ for GAD and $0.81$ for panic disorder \citep{cuijpers2016update}. A 2025 network of 375 trials gave anxiety-spectrum effects between $0.5$ and $1.0$, above $1.0$ for PTSD and specific phobia \citep{cuijpers2025unified}. In social anxiety, individual CBT ranked best of the therapies compared in a network meta-analysis of 101 trials \citep{mayowilson2014sad}, and in GAD only CBT remained significantly better than treatment as usual at 3--12 months in a network of 65 trials \citep{papola2024gad}. Any critique has to explain why CBT works at all.

\subsection{The headline shrinks with better controls and bias adjustment}
\begin{itemize}
\item \textbf{Control type.} More than 80\% of anxiety trials used waitlist controls, and the few that did not showed much smaller effects \citep{cuijpers2016update}; only 17.4\% of trials in that analysis were high quality. In the 375-trial analysis, waitlist-controlled effects exceeded $g=0.94$, while care-as-usual comparisons gave $0.22$--$1.13$ \citep{cuijpers2025unified}.
\item \textbf{Bias.} In depression, psychotherapy effects fell from $g=0.70$ to $g=0.31$ after adjusting for several sources of bias, with only 23\% of studies at low risk of bias \citep{cuijpers2019eysenck}. Of NIH-funded psychotherapy trials, 23.6\% never reached publication; including them reduced the pooled effect by about 25\% \citep{driessen2015publicationbias}. An umbrella review of 102 recent meta-analyses (3{,}782 RCTs) found most target-symptom effects small and pooled psychotherapy effects of $0.34$ against placebo or treatment as usual \citep{leichsenring2022umbrella}.
\item \textbf{Response and remission.} About half of anxiety patients respond or remit \citep{loerinc2015response,springer2018remission}.
\end{itemize}

\subsection{Does the effect shrink over time?}
This is contested, and the evidence is mixed. In depression, one meta-analysis reported a steady decline in CBT effects since its introduction \citep{johnsen2015decline}, but a reanalysis restricted to randomised comparisons judged the decline ``most likely a spurious finding'' \citep{cristea2017notfalling}, and a 409-trial meta-analysis found no reduction by publication year \citep{cuijpers2023cbtdepression}. In anxiety, a meta-regression of 49 controlled studies over 30 years found a flat trend (publication year $B=-0.008$, $p=.24$): effects ``did not increase over the last 3 decades'' \citep{hofmann2025effectsizes}. Placebo-controlled anxiety estimates have been reported at lower values in successive reviews (0.73, 0.56, 0.24), but those come from separate reviews with different case mixes, not from one meta-regression. We therefore do \emph{not} claim a decline over time; the defensible statement is that effects have not grown and that they are sensitive to control condition and bias.

\section{Where the active ingredients may lie}\label{sec:ingredients}
\subsection{Component and dismantling evidence}
The foundational dismantling trial (depression, $N=150$) found no evidence that the complete cognitive treatment outperformed behavioural activation alone or behavioural activation plus automatic-thought skills at termination or at 6 months \citep{jacobson1996component}, and a large non-inferiority trial found behavioural activation as effective as CBT (mean PHQ-9 difference $0.1$; \citealt{richards2016cobra}). In anxiety:
\begin{itemize}
\item In a 108-trial meta-analysis, cognitive therapy, exposure and their combinations were efficacious with ``no differential efficacy for any treatment components'' \citep{norton2007metaanalytic}.
\item In a full-factorial trial of internet CBT for social anxiety ($N=464$), conditions with psychoeducation and exposure were more effective, while conditions with cognitive restructuring or attention training were not superior at post-treatment; the advantage of exposure faded at follow-up \citep{sipka2025factorial}.
\item Adding cognitive restructuring to prolonged exposure did not enhance outcome in assault survivors with PTSD \citep{foa2005pecr}, nor in children with social anxiety \citep{demooij2023social}.
\item A component network meta-analysis in panic disorder (72 studies) singled out interoceptive exposure, not cognitive restructuring, as efficacy-enhancing \citep{pompoli2018dismantling}.
\item In GAD, the network meta-analysis authors concluded that pure cognitive restructuring was not a decisively active ingredient when considered alone \citep{papola2024gad}; the paper was later corrected for a duplicated cohort and an error in standard-error handling \citep{papola2024correction}.
\end{itemize}

\subsection{Evidence against the claim}
These findings are not uniform.
\begin{itemize}
\item Cognitive therapy beat exposure in social phobia in a meta-analysis of 20 direct comparisons ($N=1{,}308$), with no difference in PTSD, OCD or panic \citep{ougrin2011exposure}.
\item In a trauma trial, adding cognitive restructuring to exposure was associated with PTSD rates of 31\% versus 63--75\% without it at 6 months \citep{bryant2008cr}.
\item In-session cognitive restructuring correlates with outcome ($r=.35$), though from only four studies and correlationally \citep{ezawa2023restructuring}.
\item In a single-session trial ($N=249$), restructuring before exposure worked about as well as behavioural experiments and better than exposure without explicit processing \citep{johnco2025restructuring}.
\item The exposure--cognition split is blurred: exposure may itself work by changing the expectancy of harm \citep{hofmann2008extinction}.
\end{itemize}

\subsection{Why the component evidence is weak on both sides}
A review of 16 component studies concluded that they lack ``the statistical power nor the quality to draw any meaningful conclusion,'' because none except one could detect an effect smaller than $g=0.55$ \citep{cuijpers2017components}. Mediation evidence is correlational: cognitive change tracks symptom change however it is produced \citep{cristea2015dysfunctional,lorenzoluaces2015complicated}, cognitive mediation of symptom reduction was limited in a depression cohort \citep[median $r=.06$;][]{vittengl2014mediators}, and reviewers say no common or specific factor meets causal criteria \citep{kazdin2007mediators,cuijpers2019commonfactors}. A null component result is thus weak evidence that restructuring adds nothing, and a positive correlation is weak evidence that it matters.

\subsection{The strongest defence: practice when calm}
A defender can answer that patients do not restructure during a panic; they practise when calm so the response becomes automatic. We accept the force of this. It changes the claim, though: if rehearsal during calm periods does the work, the active ingredient is practice and habit formation, and the relevant question is the executive cost of each rehearsal. Exposure, which is central to anxiety CBT, is behavioural, and optimised exposure protocols are explicitly framed around expectancy violation and extinction \citep{craske2014maximizing}.

\section{Alternatives and their evidence}\label{sec:alternatives}
\subsection{Mindfulness and acceptance-based treatments}
The alternative that follows from the argument is to train attention and acceptance directly rather than dispute thoughts. Brewer frames anxiety as a habit loop (trigger, worry, a feeling of control) and proposes noticing, curiosity about the reward value of worry, and a ``bigger, better offer'' \citep{brewer2021unwinding}. Evidence for the specific programme is early: in a randomised trial of an app-based programme for GAD ($n=65$), anxiety scores fell 67\% versus 14\% under usual care, with mediation implicating reduced worry \citep{roy2021app}. The trial is small, self-reported, compared with usual care rather than CBT, and has the programme's developer among its authors.

Against evidence-based comparators, mindfulness and acceptance approaches show parity, not superiority. A meta-analysis found essentially no difference ($g\approx0.00$) between mindfulness-based interventions and evidence-based treatments, often CBT \citep{goldberg2018mbi}; another found no difference from traditional CBT or behavioural therapies ($g=-.07$, nine comparisons; \citealt{khoury2013mindfulness}). For acceptance and commitment therapy (ACT) versus CBT, the pooled difference was non-significant in one meta-analysis \citep{atjak2015act}, and a sceptical reading of 60 RCTs concluded that ACT is not yet well-established for any disorder \citep{ost2014act}. In 276 adults with anxiety disorders, eight weeks of mindfulness-based stress reduction was noninferior to escitalopram, with fewer discontinuations for adverse events \citep{hoge2023mbsr}. That is parity with a first-line drug, not superiority to CBT. The convergence of camps is itself informative: if CBT, behavioural activation, ACT and mindfulness are roughly equivalent, either common factors dominate or several mechanisms lead to the same place.

\subsection{Medication in ADHD}
In a network meta-analysis of 133 double-blind trials, amphetamines (standardised effect $-0.79$) and methylphenidate ($-0.49$) outperformed placebo in adults on clinician ratings, over short follow-up \citep{cortese2018adhdmeds}. The certainty varied from high or moderate to low or very low. A Cochrane review of extended-release methylphenidate found a self-rated effect of $-0.37$, judged 20 of 21 trials at high risk of bias, and noted that all trials excluded participants with psychiatric comorbidity such as anxiety or selected previous responders \citep{boesen2022ermph}. So the pivotal medication trials say little about adults with ADHD \emph{and} anxiety.

\section{CBT in adult ADHD}\label{sec:adhdcbt}
\subsection{Efficacy}
A randomised trial of 86 medication-treated adults found that CBT outperformed relaxation with educational support, with 53\% versus 23\% responders on clinician global ratings \citep{safren2010cbt}; a smaller earlier trial found 56\% versus 13\% responders for CBT added to medication \citep{safren2005cbt}. A meta-analysis of 32 studies found a between-group effect of $g=.65$ on symptoms, smaller with active controls, and neither medication status nor treatment length moderated outcome \citep{knouse2017meta}. Another found SMD $0.76$ against waitlist but $0.43$ against active controls \citep{young2020cbt}, and a 2026 meta-analysis of 14 RCTs found SMD $-0.45$ on core symptoms and $-0.43$ on executive function \citep{liu2026cbt}. A component network meta-analysis of 113 RCTs found CBT better than placebo only on clinician-reported measures \citep{ostinelli2025adultnma}, and blinded-rater effects of psychological treatments were attenuated to non-significance in an earlier synthesis that did not separate adult from child samples \citep{sonugabarke2013nonpharm}. An umbrella review judged CBT, atomoxetine and methylphenidate to have at least moderate-certainty evidence in adults, with no high-certainty long-term evidence for any intervention \citep{gosling2025umbrella}. The Cochrane review rates the CBT evidence as low quality \citep{lopez2018cochrane}.

\subsection{What the protocols contain, and where CBT sits}
The effective protocols are skills-heavy: time management, organisation and planning, with adaptive thinking as one module \citep{safren2010cbt,solanto2010metacognitive}. One review argued that learning behavioural compensatory strategies ``may be a critical active ingredient'' \citep{knouse2010current}, and the authors of one trial proposed behavioural control as the key mechanism \citep{young2017functional}. We found no dismantling study of cognitive restructuring in adult ADHD, so this is inference. In the largest multi-centre trial ($N=433$), group CBT did not beat structured clinical management, while methylphenidate beat placebo \citep{philipsen2015group}. Guidelines position CBT as an adjunct to medication \citep{kooij2019european,nice2018ng87}. Results on adding medication to CBT are mixed: no incremental benefit in a small double-blind pilot \citep{weiss2012cbt}, but greater early improvement in an unblinded trial ($N=88$), with the gap narrowing over follow-up \citep{cherkasova2020cbtmed}.

\subsection{Anxiety in ADHD}
In a meta-analysis of 13 RCTs ($n=744$), CBT-based interventions reduced anxiety in ADHD modestly (SMD $-0.29$ post-treatment, $-0.42$ at follow-up), with effects more apparent in adults; most trials delivered ADHD-focused rather than anxiety-targeted CBT \citep{wang2026cbtanxiety}. Another meta-analysis found CBT improved anxiety, particularly at follow-up, predicted by core-symptom reduction \citep{lopezpinar2020internalizing}. In children with primary anxiety in group CBT ($N=842$), neither ADHD diagnosis nor subtype predicted response or remission, which cuts against the strong form of our thesis \citep{gould2018cbtanxiety}. The adult version of that question appears untested.

\section{Does executive function moderate CBT outcome?}\label{sec:moderators}
This is the question on which the argument stands or falls. The direct evidence is thin and mostly from older adults.
\begin{itemize}
\item \textbf{Supports (at the level of the skill).} In older adults, poorer cognitive flexibility predicted poorer cognitive-restructuring quality and less distress reduction from restructuring, but \emph{not} overall CBT outcome: those with poorer flexibility ``can still benefit from standardised CBT'' \citep{johnco2014flex}.
\item \textbf{Supports (pilot).} In a pilot of CBT for GAD in older adults, participants whose executive scores remained low did not respond, while those whose scores improved responded as well as an intact-executive group \citep{mohlman2005ef}.
\item \textbf{Null or mixed.} No baseline executive effect was found in adult OCD \citep{vandborg2016memory}, child anxiety \citep{godovich2020executive} or brief psychotherapy for late-life depression \citep{schurr2025executive}; a systematic review of late-life anxiety found no replicated executive predictor \citep{kyriacou2023predictors}. In a meta-analysis of antidepressant response in late-life depression, only planning and organisation predicted outcome \citep{pimontel2016executive}.
\item \textbf{Opposite direction.} Better executive performance went with non-response to exposure-based CBT in paediatric OCD \citep{hybel2017executive}, and worse set-shifting went with better CBT response in late-life depression \citep{goodkind2016executive}.
\item \textbf{Adaptation helps.} Problem-solving therapy beat supportive therapy in older adults with executive impairment, with a larger advantage at greater impairment \citep{arean2010problem}, which points to adapting therapy rather than abandoning it.
\end{itemize}
Taken together, ``low executive function means CBT fails'' is not supported. A narrower claim is: executive capacity may limit uptake of the \emph{cognitive skill} without limiting overall response, because CBT also works through routes that do not need it. We found no adequately powered study of working-age adults with ADHD and anxiety.

\section{Non-response, dropout, relapse and harm}\label{sec:nonresponse}
\paragraph{Dropout} Dropout is the most visible failure. A meta-analysis of 115 CBT studies (20,995 participants) found 15.9\% dropout before treatment and 26.2\% during it, highest in depression and in e-therapy \citep{Fernandez2015Dropout}. CBT is not unusual here: across 669 psychotherapy studies the weighted dropout was 19.7\% and was not moderated by therapy orientation \citep{Swift2012Premature}; in 54 depression trials it was 17.5\% in psychotherapy arms \citep{Cooper2015DropoutMDD}, and in routine-care CBT for depression about 25\% \citep{Hans2013Depression}. CBT dropout is, if anything, lower than medication dropout \citep{Cuijpers2013Enduring}. Reported predictors are a weak alliance ($d=.55$; \citealp{Sharf2010Alliance}), comorbid personality disorder and minority status \citep{Cooper2015DropoutMDD}, and, in self-guided online programmes, male sex, low education and comorbid anxiety \citep{Karyotaki2015Dropout}.

\paragraph{Non-response} Even under trial conditions many patients do not respond. In a moderate-to-severe depression trial, 16-week response was 58\% and remission 40\% for cognitive therapy, with outcomes depending on therapist experience at one site \citep{DeRubeis2005CT}. Across 92 psychotherapy studies fewer than half of patients met response or remission criteria, and 48\% of care-as-usual patients also no longer met depression criteria, so the added value was about 14 percentage points \citep{Cuijpers2014Remission}. For anxiety, one review puts non-response at about a third of completers and finds that ``no predictor has been consistently supported'' \citep{Taylor2012NonResponse}; see also \citet{loerinc2015response} and \citet{springer2018remission}. Routine services do less well than trials: of 19,395 English IAPT patients, 40.3\% reliably recovered, 63.7\% reliably improved and 6.6\% reliably deteriorated \citep{Gyani2013IAPT}. The prognostic factors with support (comorbid personality disorder, chronicity, adversity) apply to all treatments (Section~\ref{sec:patients}); baseline severity did not moderate CBT versus medication in an individual-participant meta-analysis of 1,700 patients \citep{Weitz2015Severity}. In fairness, CBT-based collaborative care beat usual care in primary-care anxiety (12-month response 63.7\% versus 44.7\%; \citealp{RoyByrne2010CALM}) and adjunctive CBT roughly doubled response in treatment-resistant depression (46\% versus 22\%; \citealp{Wiles2013CoBalT}), with a benefit still present at 3--5 years \citep{Wiles2016CoBalTLong}; the Cochrane review rates the evidence for this as resting mostly on one trial \citep{Ijaz2018TRD}.

\paragraph{Relapse and durability} CBT has an enduring effect relative to stopping medication (relapse OR 2.61 favouring CBT; no significant difference from continued medication; \citealp{Cuijpers2013Enduring}; \citealp{Hollon2005Relapse}). Yet after acute-phase cognitive therapy ends, 29\% of patients relapse within one year and 54\% within two \citep{Vittengl2007Relapse}; continuation-phase CBT reduces this, especially in high-risk groups \citep{Jarrett2001Continuation}, and preventive psychological interventions cut relapse against usual care (RR 0.64; \citealp{Biesheuvel2015Recurrence}). In anxiety the follow-up record is thin: effects were not significant at 12 or more months for panic disorder, and relapse was reported in only 6 of 69 trials \citep{VanDis2020LongTerm}. In a panic trial CBT-alone response fell from 48.7\% acutely to 31.9\% six months after stopping \citep{Barlow2000Panic}.

\paragraph{Deterioration and harm} Across guided internet-CBT trials deterioration was lower with treatment than in controls (5.8\% versus 17.4\%, \citealp{Rozental2017Deterioration}; 3.36 versus 7.60, \citealp{Ebert2016Deterioration}), but low-education patients were at higher risk and did not differ from controls in the latter. Broader unwanted effects are common and poorly recorded: 5.2\% of 14,587 English therapy patients reported lasting bad effects \citep{Crawford2016Negative}; unwanted events are expected in roughly 5--20\% of patients \citep{Linden2014Rate}; therapists of 100 CBT patients reported side effects in 43 \citep{SchermulyHaupt2018Unwanted}; and a review of 85 trials found harmful events in more than 1 in 10 participants \citep{Klatte2023Adverse}. Only 21\% of 132 psychological-intervention trials monitored harms \citep{Jonsson2014Harms}, and none of 82 NIHR psychological trials mentioned an adverse event in the final report \citep{Duggan2014Adverse}. These are general to psychological therapy \citep{Lilienfeld2007Harm}, and the CBT literature has itself called for better harm detection \citep{Barlow2010Negative}. Routine outcome feedback for patients who are off track is a cheap mitigation (d = 0.23; \citealp{Delgadillo2018Feedback}; \citealp{Lambert2018Feedback}).

\paragraph{Reading} Dropout, non-response, relapse and harm are real, quantifiable and under-reported, so R4--R6 stand. They show that CBT leaves a large minority behind, not that it fails relative to alternatives, and almost none of the sources isolates the cognitive component.

\section{Patient and context factors}\label{sec:patients}
\paragraph{Comorbidity, chronicity and adversity} Comorbid personality disorder more than doubles the odds of a poor depression outcome (OR 2.16), whatever the treatment modality \citep{newtonhowes2014personality}; childhood maltreatment predicts non-response or non-remission (OR 1.43; \citealp{nanni2012maltreat}); and psychotherapy for chronic depression has a small effect against controls ($d=0.23$) and did worse than medication in direct comparisons ($d=-0.31$, attributed to dysthymia studies; \citealp{Cuijpers2010Chronic}). These are prognostic factors for all treatments, not demonstrations that CBT specifically fails. For anxiety, overall diagnostic comorbidity was generally unrelated to effect size across 148 treatment samples, with diagnosis-specific exceptions \citep{olatunji2010comorbid}. Trauma-focused CBT for PTSD in military samples shows large within-group effects ($d$ 0.78--1.10) yet about two thirds of patients retained the diagnosis, and the programmes were only marginally better than non-trauma-focused active controls \citep{steenkamp2015military}. Transdiagnostic CBT is one response to comorbidity \citep{newby2015transdiagnostic}. Baseline severity does not moderate CBT versus medication \citep{Weitz2015Severity}, and CBT still helps medication-refractory depression \citep{Wiles2016CoBalTLong}.

\paragraph{Neurodivergence, cognitive impairment, age and medical illness} The evidence on autism is thin rather than negative: ten randomized trials in autistic adults (n = 537) leave most effects very uncertain \citep{perezcastilla2026autistic}, and an earlier mixed-age meta-analysis found that effects depended on who rated them ($g=0.24$ self-report, non-significant; 0.66 informant; 0.73 clinician; \citealp{weston2016autism}). In intellectual disability, pooled effects were small ($g=0.43$, and 0.36 in randomized trials; \citealp{tapp2023id}); in dementia, CBT-family therapies had a small effect on depression (SMD $-0.23$) and an uncertain one on anxiety \citep{orgeta2022dementia}. Age itself is not a barrier: older adults recovered more often in routine IAPT data (OR 1.33; \citealp{saunders2021older}), and late-life anxiety CBT remitted 54\% of patients, though about a third relapsed \citep{johnco2026olderanxiety}. Medical comorbidity attenuates effects: CBT for chronic pain gives very small to small benefits (SMD $-0.09$ to $-0.34$; \citealp{williams2020chronicpain}), and insomnia CBT works but is weaker with medical than psychiatric comorbidity \citep{wu2015cbtimedical}. In substance use, CBT beats minimal treatment but not other specific therapies \citep{magill2019cbtsud}. We found no meta-analysis of educational attainment or IQ as a moderator. Whether adult ADHD moderates CBT outcome is addressed in Section~\ref{sec:moderators}.

\paragraph{Culture, representation and access} Culturally adapted CBT outperforms other conditions ($g=0.67$) and unadapted versions ($g=0.52$; \citealp{hall2016culturalmeta}; $d=0.32$ in \citealp{benish2011myth}, where adapting the explanation of the illness was the one significant moderator); depression trials show similar gains \citep{anik2021culturaldepression}, and CBT transports to Chinese-descent samples with modification \citep{li2023chinese}. But the trial base is narrow: across 342 depression trials, linguistic minorities were far more likely to be excluded than included and effects were almost never reported by ethnic group \citep{polo2019diversity}, in line with the broader WEIRD-sampling critique \citep{henrich2010weird}. Access shapes outcomes: in IAPT, waiting time and missed appointments predicted worse outcomes and service organisation explained 33\% of the variance in reliable improvement \citep{Clark2018IAPT}; recovery was lower in poorer areas \citep{delgadillo2016poverty}. In low-resource settings, lay-delivered CBT-family therapy works (adjusted mean difference $-4.86$ on the SSQ-14 in Zimbabwe, \citealp{chibanda2016friendship}; $-7.57$ on the BDI-II in India, \citealp{patel2017hap}; OR 0.22 for major depression in Pakistan, \citealp{rahman2008pakistan}), so there CBT fails through reach, not efficacy, with the pooled evidence rated mostly low certainty \citep{vanginneken2021primary,barbui2020lmicumbrella}.

\paragraph{Rumination and worry} Rumination predicts the onset and course of depression and overlaps with worry \citep{nolenhoeksema2008rumination}. They are also targets: rumination-focused CBT improved residual depression in a small trial with mediation by rumination change \citep{watkins2011rfcbt}, and repetitive-negative-thinking training reduced 12-month depression and anxiety diagnoses against a waitlist \citep{topper2017worryrumination}. This suggests a process that standard restructuring may miss, not a demonstrated failure.

\paragraph{Reading} R7 and R8 are supported in a limited sense: patient and context factors shift outcomes, but nearly all are shared with other treatments, and the strongest CBT-specific gaps are evidence gaps (autism, cognitive impairment, minority groups) rather than measured failures.

\section{Therapist, delivery and real-world factors}\label{sec:delivery}
\paragraph{Therapists} Therapist effects are real but modest: about 5\% of outcome variance on average (range 0.2--29\%; \citealp{johns2019therapist}), 6--7\% among low-intensity workers in IAPT \citep{firth2015pwp}, and in one large routine dataset individual therapists' recovery rates ranged from 23.5\% to 95.6\% \citep{saxon2012therapist}. Whether delivering CBT competently or adherently matters is unsettled: a meta-analysis found adherence--outcome and competence--outcome correlations near zero ($r=.02$ and $.07$; modalities pooled; \citealp{webb2010adherence}), while single CBT studies find competence effects mostly at the extremes or in subgroups \citep{branson2015competence,strunk2010competence,kuyken2008competence}. Experience alone did not improve outcomes across 6,591 patients and 170 therapists \citep{goldberg2016experience}. The alliance--outcome association is $r=.278$ across 295 studies and the same for internet therapy \citep{fluckiger2018alliance}, and sits at the therapist rather than the patient level \citep{baldwin2007alliance}; it is correlational and not CBT-specific.

\paragraph{Drift and training} Reviews argue that clinicians under-use core behavioural components and drift from ``doing'' to ``talking'' therapy \citep{waller2009drift,waller2016driftredux}; a survey found most licensed psychologists did not use exposure for PTSD \citep{becker2004exposure} (dated, US). We found no pooled prevalence of drift. Training changes behaviour only when it is active and supervised: unsupervised internet training did not change competence in a randomized trial \citep{rakovshik2016supervision}, in line with broader reviews \citep{beidas2010training,shafran2009gap}.

\paragraph{Dose, homework and schedule} Homework compliance correlates with outcome ($r=.26$, 23 studies; \citealp{mausbach2010homework}; $g=0.79$ for quantity, \citealp{kazantzis2016quality}), but patients who improve may do more homework. Improvement is curvilinear with dose, and most responders in routine care respond within 7 (low-intensity) or 14 (high-intensity) sessions \citep{robinson2020dose}; session frequency may matter more than total sessions \citep{cuijpers2013sessions}; intensified exposure scheduling gave equivalent outcomes with 32\% faster response \citep{pittig2021intensified}.

\paragraph{Does CBT survive contact with routine care?} Largely yes. Effectiveness studies for anxiety lie within the range of efficacy trials \citep{Stewart2009Anxiety}, those for depression are somewhat lower with about 25\% dropout \citep{Hans2013Depression}, and clinical representativeness did not erase effects \citep{shadish2000representative}. IAPT shows pre--post $d$ of 0.87 (depression) and 0.88 (anxiety) \citep{wakefield2021iapt}, though year-one reliable recovery was 40.3\% \citep{Gyani2013IAPT}, and a small independent diagnostic follow-up found only 9.2\% fully recovered (n = 90; single author; contested; \citealp{scott2018iapt}). Youth evidence-based practices beat usual care only modestly (mean effect 0.29; \citealp{weisz2013youthusual}).

\paragraph{Format and personalisation} Guided internet CBT is equivalent to face-to-face CBT ($g=0.05$; \citealp{carlbring2018icbt}), works in routine care \citep{etzelmueller2020routine}, and guided beats unguided only slightly at post-treatment \citep{karyotaki2021icbtnma}; unguided CBT has a small effect ($g=0.27$; \citealp{karyotaki2017selfguided}). A large pragmatic trial of supported computerised CBT in primary care found no benefit over usual care, with low uptake \citep{gilbody2015reeact}. Group and individual formats are equivalent when treatments are identical \citep{burlingame2016group}. Manualisation has no direct evidence of superiority \citep{truijens2019manual}; algorithm-matched care modestly beat stepped care (OR 1.40; \citealp{delgadillo2022stratified}; see also \citealp{firth2015stepped}).

\paragraph{Reading} R9 is partly supported. CBT works in routine care with lower and more variable yield, and the modifiable factors (supervision, dose, schedule, feedback, matching) are known, but causal evidence that poor delivery explains non-response is missing.

\section{Theory, mechanism and measurement critiques}\label{sec:mechanism}
\paragraph{Mechanisms are under-tested} \citet{kazdin2009understanding} argues that ``there is no evidence-based explanation of how or why even the most well-studied interventions produce change'', and sets standards (temporal precedence, specificity, dose--response, experiment; \citealp{kraemer2002mediators}) that most CBT mediation studies do not meet. Positive evidence exists: sudden symptom gains were preceded by cognitive change in two datasets \citep{tang1999sudden,tang2005cognitive}. It is mixed: dysfunctional-attitude change mediated symptom reduction in CBT versus interpersonal therapy, but the reverse direction fit CBT versus medication ($n=130$; \citealp{quilty2008mediational}); the CBT-side reviews call the evidence ``not yet conclusive'' \citep{driessen2010mediators} and the evidence for skill use as a mediator preliminary \citep{hundt2013skills}. Cognitive vulnerability--stress interactions predicted later depression only weakly in youth \citep{lakdawalla2007cognitive}. Defences of the cognitive model and its updates are in \citet{beck2014generic} and \citet{hofmann2013science}.

\paragraph{Is the cognitive part the active part?} In a network meta-analysis of 45 trials ($n=3{,}382$), cognitive restructuring, behavioural activation and full CBT did not differ, and all beat usual care (SMD 0.57, 0.52, 0.44) \citep{ciharova2021restructuring}. This is consistent with a behavioural reading and equally with a cognitive one; see Section~\ref{sec:ingredients} for the dismantling trials. For exposure, inhibitory learning offers an account in terms of expectancy violation rather than belief change \citep{vervliet2013extinction,weisman2018augmentation}, and a large trial found expectancy-violation-based exposure effective and optimisable \citep{pittig2021intensified}. Thought suppression can backfire \citep{wenzlaff2000suppression}, and experiential avoidance gives a third-wave counter-theory \citep{hayes1996avoidance}, though a key measure may largely track distress \citep{wolgast2014aaq} and third-wave trials were less stringent than CBT trials \citep{ost2008thirdwave}. The process-based critique of one protocol per diagnosis \citep{hayes2021processbased,hofmann2019process} comes from advocates of the alternative.

\paragraph{Common factors} Alliance ($r=.278$; \citealp{fluckiger2018alliance}), expectancy ($r=.18$; \citealp{constantino2018expectation}) and therapist effects (about 5\%; \citealp{johns2019therapist}) are replicable and modest. \citet{wampold2015common} argues they outweigh specific ingredients, and \citet{enck2019placebo} frame context effects in psychotherapy; this is contested, and it explains why trials without blinding cannot cleanly isolate CBT's specific effect.

\paragraph{Measurement} Clinician-rated scales show larger effects than self-report ($\Delta g=0.20$; \citealp{cuijpers2010selfreport}). Seven depression scales share little symptom content (mean overlap 0.36; \citealp{fried2017symptoms}); four common scales lack unidimensionality and longitudinal invariance ($n=3{,}509$; \citealp{fried2016invariance}), and sum scores hide which symptoms change \citep{fried2015sumscores}. Response is therefore partly instrument-dependent, and standardised, patient-defined outcomes are rare \citep{cuijpers2019targets,lambert2009clinical}.

\paragraph{Evidence quality} Only 11 of 115 depression trials met eight quality criteria, with $d=0.22$ versus 0.74 in the rest \citep{cuijpers2009quality}; excess significant findings and mean power of 49\% suggest small-study bias \citep{flint2014excess}; only 15\% of audited trials were prospectively registered \citep{cybulski2016registration}; and selective reporting inflated effects (SMD $-0.81$ versus $-0.54$; \citealp{miguel2021selective}). A Cochrane-method review rated all eight cognitive-therapy-versus-usual-care trials at high risk of bias \citep{jakobsen2011cttau}, though risk-of-bias tools fit psychotherapy awkwardly \citep{munder2018rob}. Fairness requires the counterweight: a multiverse analysis of 4,281 meta-analyses still gave a mean $g=0.56$, 90\% clinically relevant, though pooling all psychotherapies \citep{plessen2023multiverse}. Field-level critiques call for mechanism research and less redundant trials \citep{cuijpers2015future,tackett2019replication}.

\paragraph{Reading} R10 and the measurement and bias part of R3 are supported as limits on certainty. Nothing here shows a mechanism by which CBT specifically fails; it shows that the theory of change is under-tested and that the trial base is biased upward.

\section{The strongest objections}\label{sec:objections}
\begin{enumerate}
\item \textbf{CBT has a large evidence base.} Yes (Section~\ref{sec:evidence}). The argument concerns mechanism and moderators, not average efficacy.
\item \textbf{CBT for adult ADHD works.} It beats matched controls in medicated adults with a skills-heavy protocol; the evidence is graded low to moderate (Section~\ref{sec:adhdcbt}).
\item \textbf{Patients practise when calm.} Accepted (Section~\ref{sec:ingredients}); it relocates the active ingredient to practice.
\item \textbf{Exposure is behavioural.} Likewise; it strengthens the argument against restructuring specifically, though exposure may act via changed expectancies \citep{hofmann2008extinction}.
\item \textbf{Mindfulness is not better.} True. The argument needs only that restructuring is not uniquely necessary.
\item \textbf{Cognitive therapy sometimes wins.} Yes, in social phobia \citep{ougrin2011exposure} and in one PTSD trial \citep{bryant2008cr}.
\item \textbf{The neuroscience is about acute stress.} Correct, and the human reappraisal-under-stress findings are mixed (Section~\ref{sec:stress}).
\item \textbf{Executive function does not predict CBT outcome.} Largely true for overall outcome (Section~\ref{sec:moderators}); the argument survives only in the narrower, skill-level form.
\item \textbf{Therapists scaffold the executive load.} Plausible and untested; a therapist who generates the evidence list lowers demand and could erase a moderator effect.
\end{enumerate}
\textbf{Limitations.} Narrative review by a non-clinician; single-reviewer screening; several key studies are in older adults or children; evidence from depression stands in for anxiety on some questions (component trials, bias adjustment, decline over time).

\section{Predictions and a research agenda}\label{sec:agenda}
The argument yields predictions that could fail.
\begin{enumerate}
\item \textbf{P1 (skill-level moderation).} Baseline executive function predicts acquisition of restructuring skill more than acquisition of exposure-based skills, and predicts overall outcome only weakly.
\item \textbf{P2 (stress-dependence).} The gap between calm and in-the-moment reappraisal is larger in low-executive-function samples and at higher arousal than the moderate stressors used so far.
\item \textbf{P3 (component).} In adults with ADHD and anxiety, an exposure-and-behaviour protocol is not inferior to the same protocol plus restructuring.
\item \textbf{P4 (ADHD).} ADHD status predicts less benefit from the cognitive module specifically.
\end{enumerate}
\textbf{What would refute the argument:} an adequately powered study finding no moderation of skill acquisition by executive function; a component trial showing that restructuring adds substantial benefit over exposure; or meta-regression showing executive-function-heavy samples respond as well. The cheapest next steps are (i) an individual-participant or meta-regression analysis of existing trials (see \texttt{experiments/}), (ii) a PRISMA-style umbrella search to lift this review's method to the target journal's standard, and (iii) a power calculation for a moderator trial.

\section{Conclusion}
The ten reasons graded in Table~\ref{tab:reasons} do not support the claim that CBT fails. They support a more precise one: CBT works on average, leaves a large minority without response or with relapse, owes some of its headline effect to weak comparators, bias and measurement, rests on a theory of change that is under-tested, and is delivered with variable fidelity and unequal reach. Most of these limits are shared with other psychological treatments, and the largest CBT-specific gaps are missing studies (autism, ADHD plus anxiety, minority groups, causal tests of competence) rather than measured failures. On the executive-function argument in particular, restructuring-heavy CBT is effective on average, though less so against active controls and after bias adjustment, and its effects have not grown over three decades. What is not established is that the cognitive component is the active ingredient, or that it fails for people with low executive function. The best-supported claims are narrow: reappraisal draws on working memory; component studies, though underpowered, repeatedly fail to find an added benefit of restructuring over behavioural elements; executive capacity may limit the uptake of the cognitive skill but not overall response; and in adult ADHD the effective ingredients look behavioural. The strong, popular form of the argument, that anxiety takes the prefrontal cortex ``off-line,'' is not supported by the primary literature. Direct tests in adults with ADHD and anxiety are the missing evidence.


\section*{Declarations}
\textbf{Funding:} none. \textbf{Competing interests:} none declared. \textbf{Author note:} this is a theoretical critical review by a non-clinician; it is not medical advice. \textbf{Data and materials:} the search record, verified bibliography and analysis logs are in the public repository accompanying this paper.

\bibliographystyle{elsarticle-harv}
\bibliography{../deep-research/references}
\end{document}
